% TikZ 图：粗糙水平面上压缩弹簧
% Day 30：综合训练
\begin{tikzpicture}[scale=0.88,>=Latex,line cap=round,line join=round]
  \path[use as bounding box] (-0.35,-0.82) rectangle (7.80,2.08);
  \draw[thick] (0,0) -- (7.35,0);
  \foreach \x in {0.25,0.75,...,7.0}{
    \draw[gray] (\x,0) -- ++(0.18,-0.18);
  }
  \draw[thick] (7.35,0) -- (7.35,1.35);
  \draw[decorate,decoration={coil,aspect=0.45,segment length=4pt,amplitude=4pt},thick]
    (5.35,0.55) -- (7.35,0.55);

  % 刚接触弹簧时的位置（虚线）与最大压缩位置（实线）分开画。
  \draw[dashed,gray!80,fill=gray!6] (3.35,0.18) rectangle (4.15,0.92);
  \node[gray!80] at (3.75,0.55) {$m$};
  \draw[->,red,thick] (3.35,1.16) -- (4.15,1.16)
    node[midway,above] {$v_0$};

  \draw[fill=gray!14,thick] (4.55,0.18) rectangle (5.35,0.92);
  \node at (4.95,0.55) {$m$};
  \draw[->,orange,thick] (5.25,1.43) -- (4.35,1.43)
    node[midway,above] {$f=\mu mg$};

  \draw[dashed,blue!65] (4.15,0.08) -- (4.15,-0.35);
  \draw[dashed,blue!65] (5.35,0.08) -- (5.35,-0.35);
  \draw[<->,blue,thick] (4.15,-0.42) -- (5.35,-0.42)
    node[midway,below] {$x$};
  \node[font=\scriptsize,align=center] at (2.08,1.55)
    {虚线：刚接触弹簧};
  \node[font=\scriptsize,align=center] at (6.62,1.58)
    {最大压缩\\$v=0$};
\end{tikzpicture}
